\documentclass[11pt,aspectratio=169]{beamer}
\usetheme{Boadilla}
\usefonttheme{professionalfonts}
\setbeamertemplate{navigation symbols}{}

\usepackage[T1]{fontenc}
\usepackage{lmodern}
\usepackage{amsmath,amssymb,mathtools}
\usepackage{booktabs}
\usepackage{tabularx}
\usepackage{array}
\usepackage{graphicx}
\usepackage{xcolor}
\usepackage[authoryear,round]{natbib}
\usepackage{bibentry}
\bibliographystyle{apalike}
\setbeamertemplate{bibliography item}{}
\makeatletter
\def\hyper@natlinkstart#1{}
\def\hyper@natlinkend{}
\def\hyper@natlinkbreak#1#2{#1}
\makeatother

\graphicspath{{figs/}}

% Section headline for the audience arc.
\setbeamertemplate{headline}{%
  \hbox{}\vskip2pt\hfill\textcolor{gray!70}{\scriptsize\insertsectionhead}\hspace{1em}\vskip-6pt%
}
\setbeamertemplate{section in toc}[sections numbered]
\setbeamertemplate{frametitle}[default][left]
\setbeamerfont{framesubtitle}{size=\small}

\newcommand{\figplaceholder}[1]{\begin{center}\fbox{\begin{minipage}[c][0.32\textheight][c]{0.82\textwidth}\centering\large #1\end{minipage}}\end{center}}
\newcommand{\shortfigplaceholder}[1]{\begin{center}\fbox{\begin{minipage}[c][0.22\textheight][c]{0.82\textwidth}\centering\large #1\end{minipage}}\end{center}}
\newcommand{\compactfigplaceholder}[1]{\begin{center}\fbox{\begin{minipage}[c][0.10\textheight][c]{0.82\textwidth}\centering\normalsize #1\end{minipage}}\end{center}}
\newcommand{\R}{\mathcal{R}}
\newcommand{\E}{\mathcal{E}}
\newcommand{\dd}{\,\mathrm{d}}
\newcommand{\good}[1]{\textcolor{blue!70!black}{#1}}
\newcommand{\bad}[1]{\textcolor{red!70!black}{#1}}
\newcommand{\SAR}{speed asymmetry ratio}
% Inline slide navigation (end of sentence / bullet).
\newcommand{\app}[1]{\allowbreak\hspace{0.25em}\hyperlink{#1}{\beamergotobutton{details}}}
\newcommand{\back}[1]{\allowbreak\hspace{0.25em}\hyperlink{#1}{\beamerreturnbutton{back}}}
% Appendix / frame footer: bottom-right, not centered.
\newcommand{\navback}[1]{\par\vspace{0.35em}\noindent\hfill\hyperlink{#1}{\beamerreturnbutton{back}}}
\newcommand{\navfwd}[1]{\par\vspace{0.35em}\noindent\hfill\hyperlink{#1}{\beamergotobutton{details}}}
\newcommand{\navbackfwd}[2]{\par\vspace{0.35em}\noindent\hfill\hyperlink{#1}{\beamerreturnbutton{back}}\;\hyperlink{#2}{\beamergotobutton{details}}}

\title[Directional Congestion]{Directional Congestion}
\subtitle{Measurement, Spatial Equilibrium, and Tidal-Lane Reallocation}
\author[Duan, Feng, Gu]{Yige Duan (SJTU) $\quad$ Xiaoruo Feng (PKU) $\quad$ Yizhen Gu (PKU)}
\date[July 4, 2026]{July 4, 2026\\[0.35em]CES2026\\Southwestern University of Finance and Economics (SWUFE)\\Chengdu, China}

\begin{document}

\begin{frame}
\titlepage
\end{frame}

\begin{frame}{Roadmap}
\footnotesize\tableofcontents
\end{frame}

% =====================================================================
\section{Motivation}
% =====================================================================
\begin{frame}{Motivation}
\begin{columns}[T,onlytextwidth]
\column{0.65\textwidth}
\begin{itemize}\small
  \item \textbf{Peak-direction imbalance}
  \begin{itemize}
    \item Urban roads often face strong \bad{directional congestion}
    \item Commuting traffic is shaped by where workers live and where jobs are located
  \end{itemize}
  \item Daily two-way cost
  \begin{itemize}
    \item Return trips reverse the same flows
    \item Value-of-time costs should account for both \good{commute and return trips}
  \end{itemize}
  \item Policy margin
  \begin{itemize}
    \item \good{Tidal lanes} reallocate existing lanes across directions
    \item They may change route choices and \good{residence--workplace sorting}
  \end{itemize}
\end{itemize}
\column{0.34\textwidth}
\centering
\href{run:./A6_XiaoruoFeng_DirectionalCongestion_video.mov}{%
  \includegraphics[width=\linewidth,height=0.70\textheight,keepaspectratio]{TidalVedio/pic.jpg}}\\[3pt]
{\scriptsize \good{Click to play}}
\end{columns}
\end{frame}

% =====================================================================
\section{This Paper}
% =====================================================================
\begin{frame}{This Paper}
\small
\begin{itemize}
  \item \textbf{Measure directionality.}
  \begin{itemize}\footnotesize
    \item National pairs: median \SAR\ \good{0.878}; \good{58\%} of flow is below 0.90; more congested cities have lower \SAR\ ($r=\bad{-0.48}$, $n=99$).
    \item Within-city OD exposure is stronger (median \SAR\ \good{0.819}); road-segment asymmetry remains comparable (median \SAR\ \good{0.884}).
  \end{itemize}
  \item \textbf{Quantify what symmetry misses.}
  \begin{itemize}\footnotesize
    \item Holding fundamentals and parameters fixed, replacing directed costs with reciprocal symmetric costs changes $\bar W$ by \good{$1.10\%$} and moves \good{$3.36\%$} of residents and \good{$6.15\%$} of jobs.
  \end{itemize}
  \item \textbf{Evaluate tidal lanes under endogenous congestion GE.}
  \begin{itemize}\footnotesize
    \item Tidal top-100 reallocates one lane toward the slow direction at fixed total capacity.
    \item Welfare rises \good{$+0.87\%$} in ST and \good{$+0.32\%$} in LT.
    \item Capacity-expansion benchmarks use about \good{$1{,}100$ lane-km}; tidal lanes use \good{zero net lane-km}.
  \end{itemize}
\end{itemize}
\end{frame}

% =====================================================================
\section{Literature}
% =====================================================================
\begin{frame}{Literature}
\begin{itemize}\small
  \item \textbf{Spatial equilibrium and agglomeration}
  \begin{itemize}\footnotesize
    \item Economic geography and quantitative spatial models \citep{krugman1991,eaton2002,allen2014,ahlfeldt2015}.
    \item Transport in spatial GE with endogenous congestion \citep{allenarkolakis2022}.
  \end{itemize}
  \item \textbf{Transportation policy and congestion}
  \begin{itemize}\footnotesize
    \item Congestion theory, the fundamental law, and optimal networks \citep{vickrey1969,duranton2011,fajgelbaum2020}.
    \item Beijing evidence on congestion, pricing, and sorting \citep{barwick2024,yang2020}.
  \end{itemize}
  \item \textbf{Calibration and contribution}
  \begin{itemize}\footnotesize
    \item Beijing speed--density estimates motivate \good{$\lambda=1.1$} \citep{barwick2024}.
    \item Tidal-lane switch videos will discipline a \emph{local} lane--time response using vehicle counts and speeds; the citywide counterfactual remains conditional on the calibrated slope.
  \end{itemize}
\end{itemize}
\end{frame}

% =====================================================================
\section{Data}
% =====================================================================
\begin{frame}[label=data]{Data}
\small
\begin{columns}[T,onlytextwidth]
\column{0.38\textwidth}
\begin{itemize}
  \item \textbf{Commuting flows.}
  \begin{itemize}\footnotesize
    \item Baidu 2022 home--work commuter counts (persons, not vehicles per unit time).
  \end{itemize}
  \item \textbf{Directed speeds.}
  \begin{itemize}\footnotesize
    \item Baidu directed road-speed panel at 10-minute frequency, over one week in 2025.
  \end{itemize}
  \item \textbf{Road characteristics.}
  \begin{itemize}\footnotesize
    \item Road geometry, functional class, and lane records.
  \end{itemize}
\end{itemize}
\column{0.60\textwidth}
\centering
\includegraphics[width=\linewidth,height=0.76\textheight,keepaspectratio]{speed_eg.png}
\end{columns}
\navfwd{appdata}
\end{frame}

% =====================================================================
\section{Pattern: Directional Asymmetry}
% =====================================================================
\begin{frame}{Jobs--Housing Mismatch}
\framesubtitle{Square-grid commuting sample.}
\begin{columns}[T,onlytextwidth]
\column{0.48\textwidth}
\begin{itemize}\small
  \item \textbf{Separation.} Jobs cluster far more than residences: Duncan $D=\good{0.26}$, far above a random-allocation floor of $0.004$ (national township-level ranking: \app{appduncan}).
  \item \textbf{Direction.} AM flow is strongly radial-inward: \good{74.5\%} of core-vs-reverse trips head to the center (inbound/reverse $\good{2.9\times}$; 11/16 sectors $>0.70$). \app{appinbound}
  \item \textbf{Implication.} Separation gives peak travel a direction (inward AM, outward PM), creating the opening for tidal lanes.
\end{itemize}
\column{0.51\textwidth}
\centering\scriptsize Jobs--housing separation on the study grid\\
\includegraphics[width=\linewidth,height=0.68\textheight,keepaspectratio]{pat_jobs_housing_balance.png}
\end{columns}
\end{frame}

\begin{frame}{National Directional Asymmetry}
\framesubtitle{Township-level commuting-cost asymmetry in top 100 congested cities (Baidu 2025).}
\begin{itemize}\footnotesize
  \item \textbf{Definition.} \SAR\ $=\min(x_{ij},x_{ji})/\max(x_{ij},x_{ji})$; lower values mean stronger directional asymmetry.
  \item Flow-weighted median \SAR\ \good{0.878}; \good{58\%} of pair-flow faces \SAR\ $\leq0.90$ at AM peak.
  \item Higher Baidu commute-peak index $\Leftrightarrow$ stronger AM asymmetry ($r=\bad{-0.48}$, $n=99$).
\end{itemize}
\vspace{0.20em}
\begin{columns}[T,onlytextwidth]
\column{0.49\textwidth}
\centering
\includegraphics[width=\linewidth,height=0.42\textheight,keepaspectratio]{pat_national_ratio_hist.png}
\column{0.49\textwidth}
\centering
\includegraphics[width=\linewidth,height=0.42\textheight,keepaspectratio]{pat_congestion_am_scatter.png}
\end{columns}
\end{frame}

\begin{frame}{Observed Commuters Face Directional Friction}
\framesubtitle{Shenzhen grid-level commuting cost and commuter flow.}
\begin{itemize}\footnotesize
  \item \textbf{Unweighted prevalence.} In the grid-pair speed matrix, asymmetry is common: median \SAR\ 0.879; 57.5\% of potential pairs have \SAR\ $\leq0.90$.
  \item \textbf{Observed commuter exposure.} Counting commuters rather than grid-pair counts, \good{4.4M} observed commuters face median \SAR\ \good{0.819}; \good{71.6\%} commute on pairs with \SAR\ $\leq0.90$.
\end{itemize}
\vspace{0.25em}
\centering
\includegraphics[width=0.80\linewidth,height=0.42\textheight,keepaspectratio]{pat_sz_flow_ratio_hist.png}
\end{frame}

\begin{frame}{Segment-Level Policy Margin}
\framesubtitle{Asymmetry persists on local road segments, where capacity reallocation acts.}
\begin{itemize}\small
  \item \textbf{Road segments (AM, 3{,}670 two-way centerlines).} Median \SAR\ \good{$0.884$}; \good{$54.7\%$} have \SAR\ $\leq0.90$.
  \item Tidal lanes operate here: \good{zero net lane-km}, with capacity shifted toward the slow direction.
\end{itemize}
\vspace{0.25em}
\begin{columns}[T,onlytextwidth]
\column{0.26\textwidth}
\centering
\begin{minipage}[t][2.6em][t]{0.82\linewidth}\centering\scriptsize A. Local segment asymmetry map\end{minipage}\\[0.05em]
\includegraphics[width=0.82\linewidth,height=0.36\textheight,keepaspectratio]{pat_bj_centerline_map.png}
\column{0.26\textwidth}
\centering
\begin{minipage}[t][2.6em][t]{\linewidth}\centering\scriptsize B. Segment-level \SAR\ distribution\end{minipage}\\[0.05em]
\includegraphics[width=0.98\linewidth,height=0.40\textheight,keepaspectratio]{pat_bj_centerline_hist.png}
\column{0.44\textwidth}
\centering
\begin{minipage}[t][1.4em][t]{\linewidth}\centering\scriptsize C. How a tidal lane acts\end{minipage}\\[0.05em]
\includegraphics[width=0.98\linewidth,height=0.60\textheight,keepaspectratio]{tidal_lane_schematic.pdf}
\end{columns}
\end{frame}

% =====================================================================
\section{Model}
% =====================================================================
% Original model section, preserved from tidalraw.tex.

\begin{frame}{Model: Environment}
\begin{itemize}
  \item Locations $i,j\in N$ host $\bar{L}$ mass of residents and jobs
  \vspace{0.5em}
  \item A physical road is represented by two directed links, $(k,l)$ and $(l,k)$
  \begin{itemize}
    \item \textbf{Free-flow} times $t_{kl}^{ff}$ are taken from \textbf{night-time} speeds, so they may differ by direction
    \item A tidal-lane policy shifts one lane across directions while holding total capacity $n_{kl}+n_{lk}$ fixed
  \end{itemize}
  \vspace{0.5em}
  \item The realized edge time $t_{kl}$ is \textcolor{blue}{endogenous} and \textcolor{red}{directed}
  \begin{itemize}
    \item $\Xi_{kl}$ is model-implied commuter-passage load on link $(k,l)$, not observed vehicle flow
    \item $t_{kl}$ rises with relative load per lane; $t_{kl} \neq t_{lk}$ when directional loads differ
  \end{itemize}
\end{itemize}
\end{frame}

\begin{frame}{Individual Location Choice}
\begin{itemize}
  \item An individual $\omega$ living in $i$, working in $j$, and taking route $r\in\R_{ij}$ obtains
  \vspace{-0.5em}\begin{footnotesize}\begin{equation}
    V_{ij,r}(\omega) = \left( \frac{u_i w_j}{C_{ij,r}} \right)\varepsilon_{ij,r}(\omega),
    \qquad
    C_{ij,r} = \prod_{\ell=1}^{K_r} t_{r_{\ell-1},r_\ell}
  \vspace{-1.25em}\end{equation}\end{footnotesize}
  \begin{itemize}
    \item $u_i$ is amenity; $w_j$ is wage; $\epsilon$ is i.i.d. Frechet
    \item $C_{ij,r}$ is the route cost (product of edge times); route $r=[r_0,\dots,r_{K_r}]$, $r_0=i$, $r_{K_r}=j$
  \end{itemize}
  \vspace{0.5em}
  \item Choice probabilities are
  \vspace{-0.5em}\begin{equation}\label{eq:pi}
    \pi_{ij,r} = \frac{\left(u_i w_j/C_{ij,r}\right)^\theta}{\sum_{a}\sum_b\sum_{s\in\R_{ab}}\left(u_a w_b/C_{ab,s}\right)^\theta}
  \end{equation}
\end{itemize}
\end{frame}

\begin{frame}{Bilateral Commuting Costs and Flows}
\begin{itemize}
  \item The bilateral commuting-cost, aggregated across all routes (Frechet expectation)
  \vspace{-0.5em}\begin{equation}
    \tau_{ij} = \left[ \sum_{r\in\R_{ij}} C_{ij,r}^{-\theta} \right]^{-1/\theta}
  \end{equation}
  \item Thus, residence-workplace choice shares satisfy
  \vspace{-0.25em}\begin{equation}\label{eq:gravity}
    L_{ij} = \sum_{r\in\R_{ij}}L_{ij,r} = \tau_{ij}^{-\theta} \left(\frac{L_i^R}{\Pi_i}\right) \left(\frac{L_j^F}{P_j}\right), \,\, L_i^R=\sum_jL_{ij}, \,\, L_j^F=\sum_iL_{ij}
  \vspace{-1em}\end{equation}
  \begin{itemize}
    \item Directional congestion implies $\tau_{ij}\neq\tau_{ji}$ in general
    \item $\Pi_i$ and $P_j$ are origin and destination market-access terms
  \end{itemize}
\end{itemize}
\end{frame}

\begin{frame}{Amenity, Productivity, and Location Choices}
\begin{itemize}
  \item Agglomeration of local fundamentals
  \vspace{-0.5em}\begin{equation}\label{eq:agglom}
    w_j = A_j = \bar A_j\left(L_j^F\right)^\alpha, \qquad u_i=\bar u_i\left(L_i^R\right)^\beta
  \end{equation}
  \item Resident ($R$) and workplace ($F$) distribution across space
  \begin{align}
  (l_i^R)^{1-\beta\theta} & = \chi \sum_j\tau_{ij}^{-\theta}\,\bar u_i^{\theta}\,\bar A_j^{\theta}\,(l_j^F)^{\alpha\theta} \label{eq:clear1} \\
  (l_j^F)^{1-\alpha\theta} & = \chi \sum_i\tau_{ij}^{-\theta}\,\bar u_i^{\theta}\,\bar A_j^{\theta}\,(l_i^R)^{\beta\theta} \label{eq:clear2}
  \end{align}
\end{itemize}
\end{frame}
\note{With $\theta=6.83$, $\alpha=0.05$, and $\beta=-0.10$, the location fixed point has $1-\beta\theta\approx1.68$ on the residence side and $1-\alpha\theta\approx0.66$ on the workplace side. Negative $\beta$ is a residential dispersion force; positive $\alpha$ is workplace agglomeration. These parameters govern feedback strength and the direction of sorting rather than a stand-alone empirical result.}

\begin{frame}{Route-Specific Traversal Cost vs. Residence-Workplace Commuting Cost}
\begin{itemize}
  \item Define the directed adjacency matrix $A=[a_{kl}]$ by
  \vspace{-0.5em}\begin{equation}
    a_{kl} = t_{kl}^{-\theta}
  \end{equation}
  \item Summing over all feasible route lengths gives
  \vspace{-0.5em}\begin{equation}
    B = \sum_{K=0}^{\infty}A^K = (I-A)^{-1}
  \end{equation}
  \item The bilateral commuting-cost aggregator can be written as
  \vspace{-0.5em}\begin{equation}\label{eq:inverse}
    \tau_{ij} = B_{ij}^{-1/\theta}
  \end{equation}
\end{itemize}
\end{frame}

\begin{frame}{Feedback between Route Choice, Commuter Load, and Traversal Cost}
\begin{itemize}
  \item For a directed trip $i\to j$, define $\kappa^{kl}_{ij}$ as expected use of link $(k,l)$
  \vspace{-1em}\begin{equation}
    \kappa^{kl}_{ij} = \left( \frac{\tau_{ij}} {\tau_{ik}\,t_{kl}\,\tau_{lj}} \right)^\theta
  \vspace{-0.5em}\end{equation}
  \item The directional commuter-passage load is
  \begin{equation}\label{eq:Xi_c}
  \Xi_{kl} = \sum_i\sum_j \kappa^{kl}_{ij}L_{ij}
  \end{equation}
  \item Around the observed baseline, the congestion slope applies to relative load per lane
  \vspace{-0.5em}\begin{equation}\label{eq:c_Xi}
    \dd\log t_{kl}=\lambda\left(\dd\log\Xi_{kl}-\dd\log n_{kl}\right)
  \vspace{-0.5em}\end{equation}
  \item $\Xi_{kl}$ is measured in commuter passages per model period; the vehicle-flow bridge is an explicit proportionality assumption. \app{appflowbridge}
  \item Welfare is evaluated from the primitive Fr\'echet denominator
  \vspace{-0.5em}\begin{equation}
    \bar W
    =
    \left[
    \sum_{i,j}\tau_{ij}^{-\theta}u_i^\theta w_j^\theta
    \right]^{1/\theta},
    \qquad
    u_i=\bar u_i(L_i^R)^\beta,\quad
    w_j=\bar A_j(L_j^F)^\alpha
  \vspace{-0.8em}\end{equation}
\end{itemize}
\end{frame}

\begin{frame}{Equilibrium}
\begin{itemize}
  \item Given fundamentals $\{\bar u_i,\bar A_i,t^{\text{ff}}_{kl},n_{kl}\}$ and parameters $(\theta,\alpha,\beta,\lambda)$, an equilibrium is
  \vspace{-0.5em}\begin{equation}
    \left\{L_i^R,L_j^F,L_{ij},\tau_{ij},t_{kl},\Xi_{kl},\Pi_i,P_j\right\}
  \vspace{-0.5em}\end{equation}
  such that
  \begin{itemize}
    \item residents choose residence-workplace-route options according to \eqref{eq:pi}
    \item model-implied bilateral flows satisfy commuting gravity \eqref{eq:gravity}; residence and workplace shares clear markets \eqref{eq:clear1}--\eqref{eq:clear2}
    \item amenities and productivity satisfy \eqref{eq:agglom}
    \item bilateral costs are generated by the directed road network \eqref{eq:inverse}
    \item directed commuter load and directed traversal costs are mutually consistent \eqref{eq:Xi_c}--\eqref{eq:c_Xi}
  \end{itemize}
  \item Existence and uniqueness follow the sufficient conditions of \citet[Proposition~1]{allenarkolakis2022}. \app{appexist}
\end{itemize}
\end{frame}

% =====================================================================
\section{Estimation}\label{sec:quant}
% =====================================================================
\begin{frame}[label=estparams]{Estimation: Data and Parameters}
\small
\textbf{Data inputs}
\vspace{0.25em}

\scriptsize
\begin{tabularx}{0.92\textwidth}{@{}p{0.24\textwidth}X@{}}
\toprule
Object & Source \\
\midrule
$L_i^R,L_j^F,L_{ij}$ & Baidu 2022 commuter counts (persons per OD, not vehicle flow) \\
$t_{kl}^{\mathrm{obs}}$ & Baidu segment-level speed data: weekday AM peak average (07:00--09:00), May 19--25, 2025 \\
$t_{kl}^{ff}$ & Baidu segment-level speed data: night free-flow average (22:00--06:00, all days), May 19--25, 2025 \\
$\ell_{kl},n_{kl}$ & Road geometry and lane records \\
$(\mathcal N,\mathcal E)$ & Beijing 3 km grid and directed adjacent-grid network \\
\bottomrule
\end{tabularx}

\vspace{1.0em}
\textbf{Parameters}
\vspace{0.25em}

\scriptsize
\begin{tabularx}{0.92\textwidth}{@{}p{0.16\textwidth}p{0.15\textwidth}X@{}}
\toprule
Parameter & Value & Role / source \\
\midrule
$\theta$ & $6.83$ & cost elasticity (Ahlfeldt et al., 2015; Allen and Arkolakis, 2022) \\
$\alpha$ & $+0.05$ & production agglomeration \\
$\beta$ & $-0.10$ & residential crowding \\
$\lambda$ & $1.1$ & calibrated local speed--density elasticity (Barwick et al.\ 2024) \\
\bottomrule
\end{tabularx}
\navfwd{appparams}
\end{frame}

\begin{frame}{Estimation: Anchor and Counterfactuals}
\begin{itemize}\small
  \item \textbf{Invert fundamentals.}\\[-0.1em]
  {\footnotesize Recover $(\bar u_i,\bar A_j)$ so observed residence and employment marginals are an equilibrium at observed directed travel times.}
  \item \textbf{Anchor congestion to observed speeds.}\\[-0.1em]
  {\footnotesize Counterfactual edge times move with model-implied commuter load per lane relative to the baseline:}
  \vspace{-0.35em}\begin{equation}\label{eq:anchor}
    t_{kl}=\max\!\Big( t_{kl}^{\,\mathrm{ff}},\;\;
    t_{kl}^{\,\mathrm{obs}}\big[(\Xi_{kl}/\Xi_{kl}^{0})(n_{kl}^{0}/n_{kl})\big]^{\lambda}\Big)
  \end{equation}
  {\scriptsize $\Xi_{kl}^{0}$ is the baseline model load, not observed veh/hour. Stable person-to-vehicle conversion implies $Q_{kl}/Q_{kl}^{0}=\Xi_{kl}/\Xi_{kl}^{0}$. \app{appflowbridge}}
  \item \textbf{Reallocate lane capacity.}\\[-0.1em]
  {\footnotesize Tidal lanes hold $N_{kl}=n_{kl}+n_{lk}$ fixed and shift one lane from the fast direction to the slow direction.}
  \item \textbf{Solve in levels.}\\[-0.1em]
  {\footnotesize For each scenario, recompute $(t,\tau,\Xi,L^R,L^F,\Pi,P,\bar W_\theta)$ with fundamentals fixed; effects are differences between two equilibria.}
\end{itemize}
\end{frame}

\begin{frame}{Short Run vs.\ Long Run}
\small
\begin{itemize}
  \item \textbf{Short run (ST):} pin residence $L^R$ and employment $L^F$ at baseline; route load, edge times $t_{kl}$, and bilateral costs $\tau_{ij}$ re-equilibrate.
  \item \textbf{Long run (LT):} allow $L^R$, $L^F$, and OD shares to sort, with congestion feedback.
  \item $s_{ij}^h$ is the OD commuter share, $\tau_{ij}^h$ is bilateral commuting cost, and $\bar\tau^h=\sum_{ij}s_{ij}^h\tau_{ij}^h$ for $h\in\{0,ST,LT\}$.
\end{itemize}
\vspace{-0.5em}
{\footnotesize
  \begin{equation}
    \Delta\bar\tau^{LT}=
    \underbrace{\textstyle\sum_{ij}s_{ij}^0(\tau_{ij}^{ST}\!-\tau_{ij}^0)}_{\text{ST direct}}
    +\underbrace{\textstyle\sum_{ij}s_{ij}^0(\tau_{ij}^{LT}\!-\tau_{ij}^{ST})}_{\text{GE wedge}}
    +\underbrace{\textstyle\sum_{ij}(s_{ij}^{LT}\!-s_{ij}^0)\tau_{ij}^{LT}}_{\text{sorting}}
  \end{equation}
}
\vspace{-1.0em}
\begin{itemize}\footnotesize
  \item \textbf{Decomposition terms:}
  \begin{itemize}\scriptsize
    \item \textbf{ST direct}: baseline OD shares, baseline to short-run congestion.
    \item \textbf{GE wedge}: baseline OD shares, short-run to long-run congestion feedback.
    \item \textbf{Sorting}: OD shares move from $s^0$ to $s^{LT}$, evaluated at LT costs.
  \end{itemize}
\end{itemize}
\end{frame}

\begin{frame}[label=fundamentals]{Baseline Fundamentals (inverted)}
\begin{columns}[T,onlytextwidth]
\column{0.5\textwidth}
\centering \scriptsize Residence amenity $\bar u_i$\\
\includegraphics[width=\linewidth,height=0.52\textheight,keepaspectratio]{fund_ubar.png}
\column{0.5\textwidth}
\centering \scriptsize Workplace productivity $\bar A_j$\\
\includegraphics[width=\linewidth,height=0.52\textheight,keepaspectratio]{fund_abar.png}
\end{columns}
\vspace{0.1em}
\begin{itemize}\footnotesize
  \item $\bar u,\bar A$ are inverted to match observed residence and workplace shares at observed travel times.
  \item $\mathrm{corr}(\bar A_j,L_j^F)=\good{0.67}$ vs.\ $\mathrm{corr}(\bar u_i,L_i^R)=0.30$: jobs pinned by productivity, residence is access-driven.
  \item Treated corridors land on \bad{high-load employment bottlenecks} (TidalCandidateScore vs.\ slow-direction model load $\Xi_s$: corr.\ $\approx+0.91$ on selected top-100 pairs). \app{appfund}
\end{itemize}
\end{frame}

\begin{frame}{What Symmetry Misses}
\framesubtitle{What changes when directional congestion is treated as symmetric?}
\footnotesize
Hold the network, lanes, inverted fundamentals, parameters, solver, and welfare formula fixed; replace only reciprocal-pair $t_{kl}^{obs}$ and $t_{kl}^{ff}$ by pairwise geometric means.
\vspace{0.25em}

{\scriptsize Direction: $\Delta x=x^{\mathrm{sym}}/x^{\mathrm{dir}}-1$, with the directed baseline as the reference; dissimilarities compare $L^{\mathrm{sym}}$ to $L^{\mathrm{dir}}$.}

\vspace{0.45em}
\begin{center}
\scriptsize
\begin{tabular}{lcccc}
\toprule
 & $\Delta\bar W$ & Residence & Workplace & OD cost \\
 &  & dissim. & dissim. & $\Delta$ \\
\midrule
Symmetric vs.\ directed & \good{$+1.10\%$} & $3.36\%$ & $6.15\%$ & $+0.10\%$ \\
\bottomrule
\end{tabular}
\end{center}

\vspace{0.2em}
\begin{itemize}\footnotesize
  \item Symmetry changes location shares before any counterfactual shock.
  \item Under a hub-productivity shock, the symmetric anchored model attenuates the OD-cost response ($+1.81\%$ vs.\ $+1.93\%$).
\end{itemize}
\end{frame}

\begin{frame}{Where the Tidal Policy Lands}
\framesubtitle{Top-100 corridors by TidalCandidateScore (TCS).}
\begin{columns}[T,totalwidth=0.96\textwidth]
\column{0.42\textwidth}
\begin{itemize}\footnotesize
  \item For each reciprocal pair, label the slower AM direction $s$ and the faster direction $f$.
  \item \textbf{TCS} is the slow-minus-fast congestion-delay mass:
  \vspace{-0.35em}\[
    \mathrm{TCS}
    =\Xi_s(t_s^{obs}-t_s^{ff})-\Xi_f(t_f^{obs}-t_f^{ff}).
  \]
  \item Positive TCS means a lane shifted toward $s$ targets the direction with more model-implied commuter load times excess travel time.
  \item Red edges are the lane-receiving directions among the top-100 pairs.
\end{itemize}
\column{0.54\textwidth}
\centering
\vspace{-0.2em}
\includegraphics[width=\linewidth,height=0.78\textheight,keepaspectratio]{map_treated_overlay.png}
\end{columns}
\end{frame}

\begin{frame}{Result 1. Tidal Lanes: ST / GE / Sorting Cascade}
\begin{center}\small
\begin{tabular}{lcc}
\toprule
Channel & sec./commuter & person-min/day \\
\midrule
ST direct (fixed sorting) & \good{$-1.63$} & \good{$-2{,}400$} \\
$+$ network GE wedge & $+0.99$ & $+1{,}456$ \\
$+$ residence--workplace sorting & \bad{$+0.71$} & \bad{$+1{,}039$} \\
\midrule
\textbf{Long-run total} & \textbf{$+0.06$} & \textbf{$+95$} \\
\bottomrule
\end{tabular}
\end{center}
\vspace{0.2em}
\begin{itemize}\footnotesize
  \item \textbf{Exposure:} \good{56.4\%} of baseline OD commuter mass lies on shortest paths through treated corridors; ST direct aggregates to \good{$-2{,}400$ person-min/day}.
  \item Treated slow-direction links save \good{$20.3\%$} in ST and \good{$13.6\%$} in LT; sorting and congestion feedback absorb most of the average time gain.
\end{itemize}
\end{frame}

\begin{frame}[label=result2]{Result 2. Welfare: ST $\uparrow$, LT $\uparrow$}
\begin{center}\small
\begin{tabular}{lcc}
\toprule
Object & Short run & Long run \\
\midrule
$\bar W$ & \good{$+0.87\%$} & \good{$+0.32\%$} \\
OD route cost total & \good{$-1.63$ s} & \bad{$+0.06$ s} \\
\quad ST direct & \good{$-1.63$ s} & \good{$-1.63$ s} \\
\quad GE wedge & -- & \bad{$+0.99$ s} \\
\quad sorting wedge & -- & \bad{$+0.71$ s} \\
\bottomrule
\end{tabular}
\end{center}
\vspace{0.1em}

\vspace{0.25em}
\begin{itemize}\footnotesize
  \item Route-cost entries are seconds per baseline road commuter; negative values are time savings.
  \item Long-run sorting and congestion feedback offset the average time saving, but $\bar W$ remains positive because workers value access to the full residence--workplace opportunity set.
\end{itemize}
\end{frame}

\begin{frame}{Result 3. Long-Run Sorting}
\begin{columns}[T,onlytextwidth]
\column{0.35\textwidth}
\begin{itemize}\footnotesize
  \item Totals are fixed at $\bar L$; maps report demeaned share changes, not citywide growth.
  \item Positive cells gain share relative to baseline; black circles mark treated-corridor endpoint nodes.
  \item \textbf{Residential shares} rise near treated endpoints because lower OD costs raise residential market access.
  \item \textbf{Employment gains} are concentrated in a smaller set of red cells.
\end{itemize}
\column{0.63\textwidth}
\centering
{\scriptsize
\begin{tabular}{lcccc}
\toprule
 & \multicolumn{2}{c}{Treated endpoints} & \multicolumn{2}{c}{Rest} \\
 & $\ell^R$ & $\ell^F$ & $\ell^R$ & $\ell^F$ \\
\midrule
Mean $\Delta\log$ & \good{$+0.017$} & $-0.006$ & \bad{$-0.012$} & \bad{$-0.016$} \\
\bottomrule
\end{tabular}
}

\vspace{1em}
\begin{minipage}[t]{0.49\linewidth}
\centering \scriptsize demeaned $\Delta\log\ell^R$\\
\includegraphics[width=\linewidth,height=0.45\textheight,keepaspectratio]{sort_lR.png}
\end{minipage}\hfill
\begin{minipage}[t]{0.49\linewidth}
\centering \scriptsize demeaned $\Delta\log\ell^F$\\
\includegraphics[width=\linewidth,height=0.45\textheight,keepaspectratio]{sort_lF.png}
\end{minipage}
\end{columns}
\end{frame}

\begin{frame}[label=result4]{Result 4. Transportation Policy Comparison}
\framesubtitle{Endogenous GE under anchored congestion and long-run sorting.}
\begin{center}\scriptsize
\begin{tabular}{lcccc}
\toprule
Policy & ST $\bar W$ & LT $\bar W$ & LT route cost & net lane-km \\
 &  &  & (person-min/day) &  \\
\midrule
Tidal reallocation & $+0.87\%$ & $+0.32\%$ & $+95$ & \good{\textbf{0}} \\
Road Expansion & $+1.26\%$ & $+0.44\%$ & $+75$ & 1{,}140 \\
\bottomrule
\end{tabular}
\end{center}
\vspace{0.3em}
\begin{itemize}\footnotesize
  \item Tidal reallocation lifts LT welfare $+0.32\%$ at zero net lane-km.
  \item Delivering comparable long-run welfare by building capacity takes about $1{,}140$ lane-km: Road Expansion adds lanes where baseline commuter-load$\times$delay is largest and reaches $+0.44\%$.
  \item Expansion is a capacity benchmark, not a construction-cost estimate. \app{apppolicy}
\end{itemize}
\end{frame}

\begin{frame}[label=resultelast]{Result 5. Where to Put the Lanes}
\framesubtitle{Per-link welfare elasticity from $d\bar W^\theta=dA_{ab}\phi_a\psi_b$.}
\begin{columns}[T,onlytextwidth]
\column{0.44\textwidth}
\centering
\includegraphics[width=\linewidth,height=0.66\textheight,keepaspectratio]{map_link_welfare_elast.png}
\column{0.53\textwidth}
\scriptsize
\textbf{Selection rule} (top-100):
\begin{tabular}{lcc}
\toprule
rule & LT $\bar W$ & person-min/day \\
\midrule
\good{welfare elast.} & \good{$+0.35\%$} & $-50$ \\
TCS score & $+0.32\%$ & $+95$ \\
speed only & $+0.13\%$ & $+1430$ \\
\bottomrule
\end{tabular}
\vspace{0.5em}
\begin{itemize}\footnotesize
  \item The local score decomposes into edge gain $\times$ upstream residence access $\times$ downstream job access. \app{appelast}
  \item Model-implied load matters: speed-only rank correlation with welfare elasticity is \bad{$0.08$}; load$\times$delay correlation is $0.65$.
  \item Welfare-elasticity and TCS targeting both put capacity on the high-load spine; speed-only also selects slow but low-load edges.
\end{itemize}
\end{columns}
\end{frame}

\begin{frame}{Conclusion}
\begin{enumerate}
  \item \textbf{Directional congestion is a widespread, first-order urban feature.}
  \begin{itemize}\footnotesize
    \item Across top-100 Baidu 2025 congestion cities, flow-weighted median \SAR\ is \good{0.878}; \good{58\%} of pair-flow has \SAR\ $\leq0.90$. Within-city observed-flow exposure is stronger (median \SAR\ $0.82$), and asymmetry persists down to road segments.
  \end{itemize}
  \item \textbf{Ignoring direction mis-measures welfare.}
  \begin{itemize}\footnotesize
    \item A reciprocal-symmetric cost surface changes measured welfare by \good{$+1.10\%$} under anchored congestion GE, and moves \good{$3.36\%$} of residents and \good{$6.15\%$} of jobs.
  \end{itemize}
  \item \textbf{Tidal lanes: cost-effective welfare gains from directed reallocation.}
  \begin{itemize}\footnotesize
    \item ST route cost falls by \good{$1.63$ s/commuter} and $\bar W$ rises \good{$+0.87\%$}; LT sorting absorbs part of the time gain, but $\bar W$ remains \good{$+0.32\%$}.
    \item Supply-side policies with comparable LT welfare require about \bad{$1{,}100$ lane-km}; tidal reallocation uses \good{zero net lane-km}.
    \item Tidal targeting should combine commuter load, speed asymmetry, and access weights: upstream residence access and downstream workplace opportunity.
  \end{itemize}
\end{enumerate}
\end{frame}

% =====================================================================
\section{Next Steps}
% =====================================================================
\begin{frame}[label=nextsteps]{Next Steps}
\small
\begin{itemize}
  \item \textbf{Field measurement on implemented tidal lanes (Shenzhen).}
  \begin{itemize}\footnotesize
    \item Current videos recover direction-specific vehicle throughput around actual lane switches; the first day is descriptive because transition gaps prevent a sharp RDD.
    \item Repeated switch windows with speed or density will discipline a local policy-relevant $\lambda$, not citywide vehicle-flow levels.
    \item Use direction differences, opening and recovery, bandwidth checks, and placebo cutoffs in an RDiT design. \app{appfieldvisit}
  \end{itemize}
  \item \textbf{Identify and recover the key structural parameters.}
  \begin{itemize}\footnotesize
    \item Core parameters: $\alpha$ (workplace agglomeration), $\beta$ (residential crowding), $\theta$ (cost elasticity), $\lambda$ (speed--flow elasticity).
    \item Calibrate or estimate each as tightly as the data allow, rather than importing values from the literature.
  \end{itemize}
  \item \textbf{Two-city, two-peak, multi-grid analysis.}
  \begin{itemize}\footnotesize
    \item Run the full measurement and counterfactual pipeline for both \good{Beijing} and \good{Shenzhen}.
    \item Report \good{AM} and \good{PM} peaks separately.
    \item Check robustness across different grid resolutions.
  \end{itemize}
\end{itemize}
\end{frame}

% =====================================================================
\appendix
\section{Appendix}
% =====================================================================
\begin{frame}[label=appdata]{Appendix: Data Construction}
\begin{itemize}\small
  \item \textbf{Links.}
  \begin{itemize}\footnotesize
    \item Raw road geometry to links between adjacent nodes; drop duplicate/unusable records.
    \item Each physical road kept as symmetric infrastructure with two directed flows.
  \end{itemize}
  \item \textbf{Speeds.}
  \begin{itemize}\footnotesize
    \item $v_{klh}=\mathrm{dist}_{kl}/\mathrm{time}_{klh}$, inverse speed $m_{klh}=v_{klh}^{-1}$; AM peak selected.
  \end{itemize}
  \item \textbf{Grid and OD.}
  \begin{itemize}\footnotesize
    \item $3\times3$\,km square cells; home/work points assigned \emph{within} cell polygons and aggregated to grid OD.
    \item Restrict to the largest connected component, with a detour gate to drop implausible links.
  \end{itemize}
  \item \textbf{Lanes.}
  \begin{itemize}\footnotesize
    \item $n_{kl}$ from a corrected road-segment lane lookup; $1{,}103$ rounded unique values, $p50\approx3.37$.
  \end{itemize}
\end{itemize}
\navbackfwd{data}{appcost}
\end{frame}

\begin{frame}[label=appcost]{Appendix: Grid--Grid Travel Cost}
\begin{columns}[T,onlytextwidth]
\column{0.48\textwidth}
\begin{itemize}\footnotesize
  \item \textbf{Segment speeds.} Baidu directional crawl; AM peak $\Rightarrow$ $v_{kl}^{\mathrm{obs}}$, night window $\Rightarrow$ $v_{kl}^{ff}$.
  \item \textbf{Road network.} Directed Beijing road segments plus grid connector access points.
  \item \textbf{Pair-specific routing.} For each adjacent grid pair, shortest path on the road graph; store $t_{kl}=\ell_{kl}/v_{kl}$ on the connector.
  \item \textbf{Filters.} Drop links with routed length $>G\times$ centroid distance ($G=5$); keep the connected main sample ($N=1{,}104$).
  \item \textbf{Bilateral $\tau_{ij}$.} Model aggregator: $A_{kl}=t_{kl}^{-\theta}$, $B=(I-A)^{-1}$, $\tau_{ij}=B_{ij}^{-1/\theta}$.
\end{itemize}
\column{0.50\textwidth}
\centering\scriptsize AM observed travel time on adjacent grid links\\
\includegraphics[width=\linewidth,height=0.68\textheight,keepaspectratio]{app_edgeproj_adjacent_tau.png}
\end{columns}
\navback{data}
\end{frame}

\begin{frame}[label=appfund]{Appendix: Estimation Algorithm}
\footnotesize
\textbf{Baseline inversion}
\begin{itemize}
  \item Treat observed residence and workplace marginals as the baseline equilibrium.
  \item Recover $(\bar u_i^\theta,\bar A_j^\theta)$ by alternating updates with geometric-mean normalization.
  \item Match $(L_i^R,L_j^F)$, not the full bilateral OD matrix; OD flows are model-implied by gravity.
\end{itemize}
\vspace{0.25em}
\textbf{Counterfactual solve}
\begin{itemize}
  \item Given fundamentals and parameters, solve $(L^R,L^F)$, market access, commuter-passage load $\Xi$, and edge times jointly.
  \item Anchored congestion updates edge time around observed time and baseline model load:
  \vspace{-0.3em}
  \[
    t_{kl}=\max\{t_{kl}^{ff},\,t_{kl}^{obs}[(\Xi_{kl}/\Xi_{kl}^{0})(n_{kl}^{0}/n_{kl})]^\lambda\}.
  \]
  \vspace{-0.8em}
  \item Outer loop updates $t_{kl}$; inner loop alternates location shares and market access until convergence.
\end{itemize}
\navbackfwd{fundamentals}{appflowbridge}
\end{frame}

\begin{frame}[label=appduncan]{Appendix: Jobs--Housing Dissimilarity Benchmarks}
\framesubtitle{Duncan $D=\tfrac12\sum_g\lvert j_g/J-r_g/R\rvert$ ($j,r=$ jobs, residents per unit $g$); aspatial, $0$ balanced $\to1$ separated.}
\begin{columns}[T,onlytextwidth]
\column{0.60\textwidth}
\centering
\includegraphics[width=\linewidth,height=0.58\textheight,keepaspectratio]{app_duncan_cross_city.png}
\column{0.38\textwidth}
\begin{itemize}\footnotesize
  \item \textbf{Cross-city.} Township $D$, top-100 cities (residents $=$ origin, jobs $=$ dest mass, 2021 flows): Beijing $\good{0.20}$ ranks \good{1/100} (median $0.09$).
  \item \textbf{Caveat.} $D$ is aspatial, rising with size (corr $0.50$; $0.21$ with unit count): Beijing's lead is not a unit-count artefact, but $D$ ignores distance.
\end{itemize}
\end{columns}
\navback{data}
\end{frame}

\begin{frame}[label=appinbound]{Appendix: Radial AM Inbound--Reverse Asymmetry}
\framesubtitle{Road-mode OD; \emph{inbound} $=$ home $\geq15$ km from the job-weighted center \& work $\leq10$ km; \emph{reverse} $=$ mirror.}
\begin{columns}[T,onlytextwidth]
\column{0.62\textwidth}
\centering
\includegraphics[width=\linewidth,height=0.66\textheight,keepaspectratio]{app_inbound_sector.png}
\column{0.36\textwidth}
\begin{itemize}\footnotesize
  \item Each core-directed trip is binned into one of 16 $22.5^\circ$ sectors by \emph{home} bearing.
  \item Pooled inbound share \good{$74.5\%$}; inbound/reverse $\good{2.9\times}$; \good{11/16} sectors $>0.70$ (min $0.57$, median $0.77$, max $0.86$).
  \item Non-circular: inbound and reverse share one center, so the ratio is a genuine imbalance, not an artefact of the job-cluster construction.
\end{itemize}
\end{columns}
\navback{data}
\end{frame}

\begin{frame}{Appendix: Model-Implied Directional Asymmetry}
\framesubtitle{Baseline commuter-passage load $\Xi$: ratio of the two directions per link (model output, not vehicle counts).}
\setlength{\columnsep}{0.005\textwidth}
\begin{columns}[T,onlytextwidth]
\column{0.53\textwidth}
\centering
\includegraphics[width=\linewidth,height=0.70\textheight,keepaspectratio]{map_dir_asym.png}
\column{0.45\textwidth}
\begin{itemize}\footnotesize
  \item Color is $\log_{10}(\Xi_{\mathrm{high}}/\Xi_{\mathrm{low}})$; width is total reciprocal model load.
  \item Median two-way imbalance is \bad{$2.9\times$}; p90 is \bad{$20\times$}.
  \item High-imbalance links concentrate on the central high-load spine.
  \item Object: equilibrium commuter load $\Xi$, not observed veh/hour or observed speed.
\end{itemize}
\end{columns}
\navback{data}
\end{frame}

\begin{frame}[label=appparams]{Appendix: Parameter Interpretation}
\begin{itemize}\small
  \item $\theta$ governs sensitivity to bilateral commuting costs. Higher $\theta$ makes households and jobs respond more sharply to directed cost changes.
  \item $\alpha>0$ is workplace agglomeration in $w_j=\bar A_j(L_j^F)^\alpha$. It reinforces employment concentration around locations that gain access.
  \item $\beta<0$ is residential crowding in $u_i=\bar u_i(L_i^R)^\beta$. It penalizes residential pile-up and dampens household concentration.
  \item $\lambda$ maps relative commuter load per lane into directed edge time under the stable person-to-vehicle conversion assumption. Higher $\lambda$ raises the short-run congestion gain from reallocation and changes the long-run feedback.
  \item Sorting reflects the interaction of all four forces: accessibility changes shift OD costs, $\alpha$ pulls jobs toward stronger access, and $\beta$ limits residential concentration.
\end{itemize}
\navback{estparams}
\end{frame}

\begin{frame}[label=appsens]{Appendix: Sensitivity Summary}
\framesubtitle{Final mainline rerun, TCS top-100, primitive $\bar W$, $\lambda=1.1$.}
\begin{columns}[T,onlytextwidth]
\column{0.5\textwidth}
\centering Vary $\alpha$ ($\beta=-0.10$)\\[2pt]
\scriptsize
\begin{tabular}{cccc}
\toprule
$\alpha$ & sort (s) & person-min/day & $\bar W$ LT \\
\midrule
$-0.12$ & $-3.63$ & $-6417$ & $+0.37\%$ \\
$0$     & $-1.33$ & $-2960$ & $+0.34\%$ \\
$+0.03$ & $-0.25$ & $-1343$ & $+0.33\%$ \\
$\mathbf{+0.05}$ & $\mathbf{+0.71}$ & $\mathbf{+95}$ & $\mathbf{+0.32\%}$ \\
$+0.10$ & $+4.88$ & $+6340$ & $+0.29\%$ \\
\bottomrule
\end{tabular}
\column{0.5\textwidth}
\centering Vary $\beta$ ($\alpha=+0.05$)\\[2pt]
\scriptsize
\begin{tabular}{cccc}
\toprule
$\beta$ & sort (s) & person-min/day & $\bar W$ LT \\
\midrule
$\mathbf{-0.10}$ & $\mathbf{+0.71}$ & $\mathbf{+95}$ & $\mathbf{+0.32\%}$ \\
$0$     & $-1.57$ & $-3173$ & $+0.31\%$ \\
$+0.05$ & $-4.64$ & $-7579$ & $+0.28\%$ \\
$+0.10$ & \multicolumn{3}{c}{\emph{not converged}} \\
\bottomrule
\end{tabular}
\end{columns}
\vspace{0.25em}
\begin{itemize}\footnotesize
  \item $\bar W$ LT remains positive in every converged $\alpha$--$\beta$ cell; route-time statistics move more because sorting changes OD weights.
  \item Detailed cells, $\theta$ grid, and congestion-band checks\app{appsensdetail}\app{appsens_theta}.
\end{itemize}
\navback{data}
\end{frame}

\begin{frame}[label=appsensdetail]{Appendix: Sensitivity Details}
\scriptsize
\begin{columns}[T,onlytextwidth]
\column{0.40\textwidth}
\begin{flushright}
\setlength{\tabcolsep}{3.5pt}
\begin{tabular}{ccccc}
\toprule
$\alpha$ & $\beta$ & sort (s) & person-min/day & $\bar W$ LT \\
\midrule
$-0.12$ & $-0.10$ & $-3.63$ & $-6{,}417$ & $+0.37\%$ \\
$0.00$ & $-0.10$ & $-1.33$ & $-2{,}960$ & $+0.34\%$ \\
$0.03$ & $-0.10$ & $-0.25$ & $-1{,}343$ & $+0.33\%$ \\
$\mathbf{0.05}$ & $\mathbf{-0.10}$ & $\mathbf{+0.71}$ & $\mathbf{+95}$ & $\mathbf{+0.32\%}$ \\
$0.10$ & $-0.10$ & $+4.88$ & $+6{,}340$ & $+0.29\%$ \\
$0.05$ & $0.00$ & $-1.57$ & $-3{,}173$ & $+0.31\%$ \\
$0.05$ & $0.05$ & $-4.64$ & $-7{,}579$ & $+0.28\%$ \\
$0.05$ & $0.10$ & \multicolumn{3}{c}{\emph{not converged}} \\
\bottomrule
\end{tabular}
\end{flushright}
\column{0.58\textwidth}
\begin{itemize}\footnotesize
  \item ST direct effect is stable: about $-1.63$ s/commuter and $\bar W$ ST about $+0.87\%$.
  \item $\theta$ grid: each value re-inverts fundamentals and GE; LT $\bar W$ gains decline with $\theta$ while signs stay positive. \app{appsens_theta}
  \item Congestion band: at $\lambda=0.57$, top-100 LT person-min/day is about $-1{,}168$ and $\bar W$ LT is $-0.078\%$; this is retained as a legacy lower-slope robustness, not the mainline.
  \item Selection robustness: the asymmetry-ratio rule is weaker than TCS selection and can flip LT $\bar W$ negative.
\end{itemize}
\end{columns}
\navback{appsens}
\end{frame}

\begin{frame}[label=appsens_theta]{Appendix: Theta Sensitivity}
\framesubtitle{TCS top-100; each $\theta$ re-inverts fundamentals and re-solves baseline + tidal ST/LT.}
\scriptsize
\begin{center}
\renewcommand{\arraystretch}{1.3}
\setlength{\tabcolsep}{12pt}
\begin{tabular}{ccccc}
\toprule
$\theta$ & ST $\bar W$ & LT $\bar W$ & res.\ dissim & job dissim \\
\midrule
$3$ & $+1.99\%$ & $+0.75\%$ & $1.11$ & $1.55$ \\
$5$ & $+1.19\%$ & $+0.45\%$ & $0.97$ & $1.67$ \\
$\mathbf{6.83}$ & $\mathbf{+0.87\%}$ & $\mathbf{+0.32\%}$ & $\mathbf{0.86}$ & $\mathbf{1.79}$ \\
$8$ & $+0.74\%$ & $+0.27\%$ & $0.80$ & $1.89$ \\
$10$ & $+0.59\%$ & $+0.21\%$ & $0.70$ & $2.07$ \\
\bottomrule
\end{tabular}
\end{center}
\vspace{0.4em}
\begin{itemize}\footnotesize
  \item LT $\bar W$ stays positive across the grid; magnitudes fall as $\theta$ rises (stronger cost sensitivity).
  \item Resident dissim falls and job dissim rises with $\theta$; baseline $\bar W$ levels are not comparable across $\theta$.
\end{itemize}
\note{Mainline $\theta=6.83$ matches the tidal diagnostics row; baseline $\bar W$ levels are not comparable across $\theta$.}
\navback{appsens}
\end{frame}

\begin{frame}[label=apppolicy]{Appendix: Policy Definitions and Selected Effects}
\scriptsize
\begin{center}
\begin{tabular}{p{0.21\textwidth}p{0.43\textwidth}ccc}
\toprule
Policy & Capacity rule & ST $\bar W$ & LT $\bar W$ & lane-km \\
\midrule
Tidal reallocation & Use the TCS top-100 corridors; add one lane in the slow direction and remove one lane from its reciprocal fast direction. & $+0.87\%$ & $+0.32\%$ & 0 \\
Road Expansion & Re-rank reciprocal corridors by baseline commuter load $\times$ excess travel time; for the top 100, add one lane to each direction ($n_{kl}$ and $n_{lk}$ both rise by 1). & $+1.26\%$ & $+0.44\%$ & 1{,}140 \\
Uniform expansion & Scale all links by $+2.4\%$ so added lane-km matches the targeted-expansion budget; no lanes removed. & $+0.21\%$ & $+0.10\%$ & 1{,}092 \\
\bottomrule
\end{tabular}
\end{center}
\vspace{0.2em}
\begin{itemize}\footnotesize
  \item ST pins residence and workplace shares at baseline; LT re-solves residence--workplace sorting.
  \item Road Expansion asks where an engineer would add capacity using model-implied commuter load and delay only; uniform spreads the same lane-km budget across the whole network (targeting versus spreading).
\end{itemize}
\navback{result4}
\end{frame}

\begin{frame}[label=appelast]{Appendix: Welfare-Elasticity Score}
\scriptsize
\begin{itemize}
  \item Define the directed-link accessibility matrix by $A_{ab}=t_{ab}^{-\theta}$ and the network multiplier by $B=(I-A)^{-1}$.
  \item A small improvement on one directed edge uses the matrix-inverse differential:
  \[
    dB=B(dA)B.
  \]
  \item If only edge $a\to b$ changes, $dA$ has one nonzero entry, so the welfare-index numerator has the local derivative
  \[
    d\bar W^\theta=dA_{ab}\phi_a\psi_b.
  \]
  \item $dA_{ab}$ is the direct gain from one more lane; $\phi_a$ is upstream residence-access weight reaching node $a$; $\psi_b$ is downstream workplace-opportunity weight from node $b$.
  \item The score is a local ranking device. The reported LT $\bar W$ values still come from re-solving the endogenous GE counterfactual.
\end{itemize}
\navback{resultelast}
\end{frame}

\begin{frame}[label=appflowbridge]{Appendix: Commuter Load and Vehicle Flow}
\small
\begin{itemize}
  \item \textbf{Model object.} $\Xi_{kl}$ is OD-implied commuter passages on link $(k,l)$ per representative commute period. It is not veh/hour.
  \item \textbf{Vehicle-flow bridge.} Let $Q_{kl}$ denote passenger-car-equivalent flow and write
  \vspace{-0.25em}
  \[
    Q_{kl}=\frac{\rho_{kl}}{H}\Xi_{kl},
    \qquad
    \frac{Q_{kl}}{Q_{kl}^{0}}
    =\frac{\Xi_{kl}}{\Xi_{kl}^{0}}
    \quad\text{if }\rho_{kl}\text{ and }H\text{ are stable.}
  \]
  \vspace{-0.65em}
  {\footnotesize $\rho_{kl}$ absorbs mode composition, vehicle occupancy, and passenger-car-equivalent weights; $H$ is the peak-period duration.}
  \item \textbf{Interpretation.} The anchored counterfactual identifies relative load changes conditional on stable conversion. It does not recover citywide vehicle-flow levels.
  \item \textbf{Field discipline.} Switch-time video can estimate a local lane--time response from vehicle throughput plus speed or density. This local moment calibrates or bounds $\lambda$; it does not replace $\Xi$ with observed traffic.
\end{itemize}
\navback{estparams}
\end{frame}

\begin{frame}[label=appfieldvisit]{Appendix: Field Visit}
\framesubtitle{Current videos measure directional vehicle throughput; local speed or density response remains the identifying outcome.}
\vspace{-0.15em}
\setlength{\columnsep}{0.014\textwidth}
\begin{columns}[T,onlytextwidth]
\column{0.317\textwidth}
\centering
\includegraphics[width=\linewidth,height=0.305\textheight,keepaspectratio]{fieldvisit_crop_1.jpg}\par\vspace{0.24em}
\includegraphics[width=\linewidth,height=0.305\textheight,keepaspectratio]{fieldvisit_crop_2.jpg}
\column{0.317\textwidth}
\centering
\includegraphics[width=\linewidth,height=0.305\textheight,keepaspectratio]{fieldvisit_crop_5.jpg}\par\vspace{0.24em}
\includegraphics[width=\linewidth,height=0.305\textheight,keepaspectratio]{fieldvisit_crop_4.jpg}
\column{0.317\textwidth}
\centering
\includegraphics[width=\linewidth,height=0.305\textheight,keepaspectratio]{fieldvisit_crop_6.jpg}\par\vspace{0.24em}
\includegraphics[width=\linewidth,height=0.305\textheight,keepaspectratio]{fieldvisit_crop_7.jpg}
\end{columns}
\navbackfwd{nextsteps}{appflowbridge}
\end{frame}

\begin{frame}[label=appexist]{Appendix: Existence and Uniqueness}
\framesubtitle{Urban commuting model; sufficient conditions from \citet[Proposition~1]{allenarkolakis2022}.}
\small
\begin{itemize}
  \item \textbf{Existence.} On a strongly connected network a strictly positive equilibrium always exists, for any $\alpha,\beta$ and any $\theta>0$, $\lambda\ge0$ (Brouwer fixed point on the transformed operator).
  \item \textbf{Uniqueness (sufficient), for $\alpha,\beta\in[-1,1]$:}
  \vspace{-0.15em}\[
    \alpha\le\tfrac12\!\left(\tfrac1\theta-\lambda\right)
    \qquad\text{and}\qquad
    \beta\le\tfrac12\!\left(\tfrac1\theta-\lambda\right).
  \]
  \item \textbf{Principle.} The proof bounds the equilibrium operator so that its elasticity map has spectral radius at most one \citep{allenarkolakisli2020,karlin1967}.
  \item \textbf{Intuition.} $\alpha$ (productivity) and $\beta$ (amenity) are agglomeration spillovers; uniqueness requires dispersion to dominate concentration, and stronger congestion $\lambda$ lowers the admissible $\alpha,\beta$.
  \item The condition is sufficient, not necessary. At $(\alpha,\beta)=(0.05,\,0.10)$ both spillovers turn agglomerative, so the iteration fails to converge.
\end{itemize}
\navback{data}
\end{frame}

\nobibliography{references}

\end{document}
